\section{Hello}

“我能吞下玻璃而不伤身体”是中文字体测试中常用的示例文字，其来源是一个语言学项目，而非真实行为。

\subsection{背景来源}
这句话的原文是英文“I can eat glass, it doesn't hurt me.”，由哈佛大学学生Ethan在1997年发起的“I Can Eat Glass”项目创造。该项目旨在收集不同语言中这句话的翻译，用于展示旅行者在陌生语言环境中说出意想不到的句子，从而引起当地人的注意与尊重。项目最终收录超150种语言版本，涵盖自然语言、人造语言与计算机语言，迅速成为经典网络模因。

\subsection{为何成为中文字体测试文本}
这句话能广泛用于字体预览、编码校验与排版测试，核心原因有三点：
\begin{enumerate}
    \item 长度适中：语句简短完整，适配各类界面与预览窗口，不占过多空间。
    \item 覆盖常用字形：包含左右结构、上下结构、独体字等典型汉字，可直观检验笔画清晰度、间距与字形完整性。
    \item 编码兼容性强：无生僻字，在GBK、UTF-8等主流编码下稳定显示，是排查乱码、字体回退的理想样本。
\end{enumerate}

它最早被哥伦比亚大学UTF-8测试项目采用，后被Ubuntu、GNOME等开源系统与工具沿用，成为中文技术圈辨识度极高的测试语料。

这句话仅为趣味语言学创作与技术测试文本，现实中吞食玻璃会严重划伤口腔、食道与肠胃，危及生命，绝对不可模仿。


\lipsum[1-5]

